\subsection{The pentagon's dihedral symmetry $D_5$}

The KCBS pentagon (\cref{ch:relational-physics-contextuality}) puts its five measurements on the
vertices of a 5-cycle $C_5$, exclusivity running along the edges. This subsection records that
cycle's symmetry group, built ahead of the KCBS scenario itself so the possibilistic layer can
later cite it directly.

\begin{proposition}[The carrier $D_5$ and its action on the pentagon]\label{thm:d5-carrier}\lean{Obsmetrologia.QuantumContextuality.D5, Obsmetrologia.QuantumContextuality.D5.act}\leanok\uses{thm:d8-carrier, thm:measurement-scenario}
\texttt{Quantum.D5} is the ten-element semidirect-product pair $(\mathrm{rot} : \mathrm{Fin}\,5,
\mathrm{flip} : \mathrm{Bool})$: five rotations and five reflections of the pentagon. Its action
\texttt{D5.act} sends a rotation to $i \mapsto \mathrm{rot} + i$ and a reflection to $i \mapsto
\mathrm{rot} - i$ on the vertex set $\mathrm{Fin}\,5$ --- the odd-order sibling of the $D_4$
construction on $\mathrm{Charge}$ (\cref{thm:d8-carrier}), here acting on the cycle that carries
the KCBS scenario (\cref{thm:measurement-scenario}) rather than on the bilattice.
\end{proposition}
\begin{theorem}[Faithful and a homomorphism]\label{thm:d5-faithful-hom}\lean{Obsmetrologia.QuantumContextuality.D5.faithful, Obsmetrologia.QuantumContextuality.D5.act_mul}\leanok\uses{thm:d5-carrier}
\texttt{D5.faithful} proves distinct elements act as distinct permutations of the pentagon;
\texttt{D5.act\_mul} proves the semidirect-product multiplication realizes composition of the
underlying $\mathrm{Fin}\,5$ maps.
\end{theorem}
\begin{theorem}[$D_5$ is the dihedral group of order 10]\label{thm:d5-is-dihedral}\lean{Obsmetrologia.QuantumContextuality.D5.r_pow_five, Obsmetrologia.QuantumContextuality.D5.s_sq, Obsmetrologia.QuantumContextuality.D5.s_r_s, Obsmetrologia.QuantumContextuality.D5.non_abelian, Obsmetrologia.QuantumContextuality.D5.card_eq_ten}\leanok\uses{thm:d5-faithful-hom}
With $r$ the order-5 rotation and $s$ a reflection, \texttt{D5.r\_pow\_five} / \texttt{D5.s\_sq} /
\texttt{D5.s\_r\_s} give the presentation $\langle r, s \mid r^5 = s^2 = 1,\, srs = r^{-1}\rangle$;
\texttt{D5.non\_abelian} rules out an abelian group of the same order; \texttt{D5.card\_eq\_ten}
fixes the carrier at exactly ten elements. Group associativity and two-sided inverses
(\texttt{D5.mul\_assoc}, \texttt{D5.inv\_mul}, \texttt{D5.mul\_inv}) complete the group laws.
\end{theorem}
\begin{theorem}[$D_5$ acts by automorphisms of $C_5$]\label{thm:d5-preserves-c5}\lean{Obsmetrologia.QuantumContextuality.D5.adj, Obsmetrologia.QuantumContextuality.D5.act_preserves_adj}\leanok\uses{thm:d5-is-dihedral}
\texttt{D5.adj} is the pentagon's cyclic adjacency, the $C_5$ exclusivity graph's edge relation;
\texttt{D5.act\_preserves\_adj} proves every element of $D_5$ preserves it --- $D_5$ acts by graph
automorphisms of $C_5$. This does not yet prove $D_5$ is the \emph{full} automorphism group of
$C_5$ (that every adjacency-preserving permutation of $\mathrm{Fin}\,5$ is one of these ten); that
stronger claim is frontier.
\end{theorem}
